% --- Новий підрозділ: вставити у §2 "Аналіз існуючих підходів" після §2.2 ---

\subsection{Записи архітектурних рішень}
\label{sec:related-adr}

Практика формального документування архітектурних рішень має тривалу історію в програмній інженерії.
Формат Architecture Decision Record (ADR), запропонований Найґардом~\citep{nygard2011adr}, визначає легковагий шаблон у форматі Markdown --- назва, статус, контекст, рішення, наслідки --- що зберігається безпосередньо разом із кодовою базою.
Стандарт ISO/IEC/IEEE~42010~\citep{iso42010} формалізує опис архітектури на організаційному рівні, надаючи нормативну рамку для документування архітектурних точок зору та відповідностей.
Y"=формат Цьорнера~\citep{zorner2015softwarearchitekturen} структурує рішення у вигляді шаблонного висловлювання: <<В контексті [ситуація], зіткнувшись із [проблемою], ми вирішили [варіант], щоб досягти [якість], приймаючи [компроміс]>>.
\citet{zimmermann2015architectural} розширює практику управління ADR до міжпроєктного керівництва, додаючи моделювання простору проблем та систематичне відстеження залежностей між рішеннями.

ADR фіксують \emph{провенанс рішень} --- саме той контент, що відсутній у витягуванні code-RAG.
Проте механізм витягування ADR залишається файловим: рішення знаходять переглядом директорії за назвою та датою, а не семантичним запитом.
Відсутні індекс вбудовувань (embedding index), гібридне витягування, перерангування за релевантністю до поточної задачі.
Крім того, ведення ADR потребує ручного авторства: у високошвидкісному робочому процесі з продуктивністю 14,7~PR/день кураторське навантаження на підтримку повного реєстру ADR стає заборонно високим.

За військовою аналогією, ADR подібні до журналу рішень командира --- цінний артефакт, але пошук у ньому здійснюється за датою та підрозділом, а не за семантичною релевантністю до поточної оперативної ситуації.

% --- Доповнення §2.2 "Витягування для програмних агентів": ---
% --- вставити після згадки RepoCoder ---

CodeRAG-Bench~\citep{wang2024coderagbench} забезпечує систематичний бенчмарк для витягування з доповненою генерацією коду, охоплюючи документацію, довідники API та приклади коду як джерела витягування.
SWE-bench~\citep{jimenez2024swebench} та SWE-agent~\citep{yang2024sweagent} оцінюють агентів на реальних задачах з GitHub (issues), демонструючи, що витягування з урахуванням структури репозиторію та контексту задачі підвищує частку успішних розв'язків.

Усі зазначені системи витягують \emph{текст коду та ідентифікатори}.
Вони здатні відповісти на запит <<де визначено цю функцію?>> або <<що робить цей модуль?>>, але не на запит <<чому обрано саме такий підхід замість альтернативного?>> або <<який принцип конституції забезпечує цей валідатор?>>.
Контекст рішень --- обґрунтування, відхилені альтернативи, перехресні обмеження --- відсутній у самому коді і, відповідно, відсутній у витягуванні code-RAG.
